%%%%%%%%%%%%%%%%%%%%%%%%%%%%%%%%%%%%%%%%%%%%%%%%%%%%%%%%%%%%%%%%%%%%%%%%%%
%   This is frontpage.tex file needed for the mhar-vell.cls file.  You   %
%  have to  put this file in the same directory with your thesis files.  %
%                Written by M. Imran 2001/06/18                          % 
%                 No Copyright for this file                             % 
%                 Save your time and enjoy it                            % 
%                                                                        % 
%%%%%%%%%%%%%%%%%%%%%%%%%%%%%%%%%%%%%%%%%%%%%%%%%%%%%%%%%%%%%%%%%%%%%%%%%%%
%%%%%%%%%%%%%%%%%%%%%%%%%%%%%%%%%%%%%%%%%%%%%%%%%%%%%%%%%%%%%%%%%%%%%%%%%%%
%%%%%%%%%%%%%%%%           The title page           %%%%%%%%%%%%%%%%%%%%%%%  
%%%%%%%%%%%%%%%%%%%%%%%%%%%%%%%%%%%%%%%%%%%%%%%%%%%%%%%%%%%%%%%%%%%%%%%%%%%
\pagenumbering{roman}
%\pagenumbering{arabic}

\setcounter{page}{1}

\newpage

\thispagestyle{empty}

\begin{center}
  \vspace*{0.8cm}
  {\Large \bf Proposta de projeto de pesquisa do programa PRH 55 da ANP com ênfase no setor de petróleo, gás e biocombustíveis, como parte do processo de formação científica.}

  \vspace*{2.5cm}
 % {\LARGE\bf Marco Antônio dos Reis}
  %\vspace*{3cm}
  %\vfill

  {\Large Robô de inspeção com capacidade de realizar operações semi-autônomas em estruturas subaquáticas utilizando ativos nacionais e de baixo custo.}
  %\\
           \vspace*{2.5cm}
  
  {\bf ---------------------------------------\\
  Luiz Paulo de Souza Santos}
  
  \vspace*{1.5cm}
  {\bf --------------------------------------------------\\
  Marco Reis, M.Sc. (Orientador)}
  
   \vspace*{1.5cm}
  {\bf ------------------------------------------------------------------\\
  Geovani Mimoso , M.Sc. (Co-orientador)}
  
  %{\LARGE\bf Marco Antônio dos Reis}
  % Put your university logo here if you wish.
   \begin{center}
   %\includegraphics{DU_2-col_sml.eps}
   \end{center}
	\vspace*{- 0.3cm}

  {Programa de Formação de Recursos Humanos da ANP – PRH 55\\  
          [-0.3cm] Linha de pesquisa em sistemas computacionais\\
          [-3mm] Senai Cimatec\\
          [-3mm] Salvador - Bahia\\
          [-3mm] março de 2016}
\end{center}

%%%%%%%%%%%%%%%%%%%%%%%%%%%%%%%%%%%%%%%%%%%%%%%%%%%%%%%%%%%%%%%%%%%%%%%%%%%
%%%%%%%%%%%%%%%% The dedication page, of you have one  %%%%%%%%%%%%%%%%%%%%  
%%%%%%%%%%%%%%%%%%%%%%%%%%%%%%%%%%%%%%%%%%%%%%%%%%%%%%%%%%%%%%%%%%%%%%%%%%%
\newpage
\thispagestyle{empty}
\begin{center}
 \vspace*{2cm}
 % \textit{\LARGE {Dedicated to}}\\ 
 %Someone here
\end{center}


%%%%%%%%%%%%%%%%%%%%%%%%%%%%%%%%%%%%%%%%%%%%%%%%%%%%%%%%%%%%%%%%%%%%%%%%%%%
%%%%%%%%%%%%%%%%%%           The abstract page         %%%%%%%%%%%%%%%%%%%%  
%%%%%%%%%%%%%%%%%%%%%%%%%%%%%%%%%%%%%%%%%%%%%%%%%%%%%%%%%%%%%%%%%%%%%%%%%%%
\newpage
\thispagestyle{empty}
\addcontentsline{toc}{chapter}{\numberline{}Resumo}
\begin{center}
  \textbf{\Large Robô de inspeção com capacidade de realizar operações semi-autônomas em estruturas subaquáticas utilizando ativos nacionais e de baixo custo.
  }

  \vspace*{1cm}
  \textbf{\large Luiz Paulo de Souza Santos}

  \vspace*{2cm}
 % {\large Submetido para proposta de projeto ao processo de admissão no Doutorado de Mecatrônica\\ janeiro de 2016}

  %\vspace*{2cm}
  \textbf{\large Resumo}
\end{center}
%Durante vários anos, a inspeção de estruturas submarinas tem sido feito regularmente por ROVs ou mergulhadores em alta profundidade. Além de ser uma forma onerosa de inspeção, em alguns casos, devido às condições meteorológicas e outros fatores que dificultam a inspeção. Como alternativa para o uso de mergulhadores, a inspeção por veículos remotamente controlados é apresentada como uma inspeção limitada, devido ao tempo despendido para a realização de tarefas simples e repetidas, e que muitas vezes não são reprodutivas e não alcançam um bom desempenho quanto a repetibilidade.
%Desta forma uma inspeção por ROVs com operações semi-autônomas torna-se uma solução vantajosa quanto à redução dos custos, riscos relacionados aos recursos e aumento de produtividade. 



Desde o surgimento da extração de petróleo e gás em ambientes subaquáticos há a necessidade de inspecionar as estruturas submersas devido a grande salinidade do mar, a movimentação da plataforma causadas pelas ondas e ao tempo de utilização das plataformas. Estas inspeções são realizadas atualmente por mergulhadores e ROV's  (Remotely Operated Vehicle) no intuito de evitar a corrosão, trincas ou rompimento das tubulações e/ou cabos que ligam a plataforma ao inferior do oceano. As altas profundidades proporcionam uma variação muito grande de pressão e temperatura as quais um ser humano não deve ser exposto, diante disso os ROV's tem uma grande vantagem, entretanto algumas limitações como os custos operacionais, e condições meteorológicas são fatores cruciais para o desenvolvimento dessas vistorias pois os processos atuais dependem de um navio de apoio offshore e mão de obra altamente especializadas.
A rotina de manutenção nas atividades em plataformas é repetitiva pois essas estruturas são fixas e quanto mais velha é a instalação mais rotinas devem ser realizadas o que faz uma aplicação local e semi autônoma de um Veículo operado remotamente um ótimo método de redução de custos e um aumento na confiabilidade das instalações.

O projeto será realizado mediante um ambiente de simulação utilizando o framework ROS e o software GAZEBO. Com isso espera-se apresentar uma alternativa à metodologia de inspeção em estruturas submarinas utilizada hoje em dia. A continuidade da pesquisa no uso de uma solução mecânica já desenvolvida é algo vital para este projeto, logo será utilizado o OpenRov, um sistema mecânico já utilizado em pesquisas anteriores para inspeções em estruturas submarinas.

Este projeto visa desenvolver, em dois anos, um supervisório capaz de operar o OpenRov  com uma interface operacional de fácil compreensão e utilização tendo algumas operações realizadas pelo ROV sendo executadas automaticamente. Será desenvolvido em duas fases consecutivas de 1 anos cada. Serão desenvolvidas as seguintes atividades:




 \textbf{\large 1ª Fase.}
 
\begin{itemize}

	\item Montagem do Open ROV.
	\item Pesquisar configuração padrão do Open ROV.
	\item Pesquisar linguagens usadas no Open ROV.
	\item Pesquisar padrões de comunicação do Open ROV.
	\item Desenvolver um algorítimo no ROS para controlar o ROV de um computador.
	\item Desenvolver um algorítimo no ROS que recebe as informações adquiridas pelo sensores do OpenRov.
\end{itemize}

 \textbf{\large 2ª Fase.}
 
\begin{itemize}

	\item Desenvolver um algorítimo no ROS que controle o OpenRov de maneira semi-autônoma de acordo com as informações adquiridas.
    \item Desenvolver a simulação no software GAZEBO.
	\item Desenvolver o TCC.

\end{itemize}


\vspace*{1cm}
\noindent \textbf{Palavras chaves: }robô de inspeção, inspeção subaquática.


%%%%%%%%%%%%%%%%%%%%%%%%%%%%%%%%%%%%%%%%%%%%%%%%%%%%%%%%%%%%%%%%%%%%%%%%%%%
%%%%%%%%%%%%%%%%%%          The declaration page         %%%%%%%%%%%%%%%%%%  
%%%%%%%%%%%%%%%%%%%%%%%%%%%%%%%%%%%%%%%%%%%%%%%%%%%%%%%%%%%%%%%%%%%%%%%%%%%
\chapter*{Declaração}
\addcontentsline{toc}{chapter}{\numberline{}Declaração}
O projeto nesta proposta é baseado na pesquisa realizada pelo Brazilian Institute of Robotics - BIR / Senai Cimatec. A pesquisa versa na área da inspeção de estruturas subaquáticas. Referências contidas nela são estritamente de responsabilidade do autor.\\
\vspace*{8cm}
\vfill
\noindent \textbf{Copyright\copyright\; 2016 by Luiz Paulo de Souza Santos}.\\
``Os direitos de autoria deste projeto recai sobre o autor. Citações não devem ser publicadas sem o consentimento prévio por escrito do autor e informações oriundas da mesma devem ser reconhecidas e referenciadas''.



%%%%%%%%%%%%%%%%%%%%%%%%%%%%%%%%%%%%%%%%%%%%%%%%%%%%%%%%%%%%%%%%%%%%%%%%%%%
%%%%%%%%%%%%%%%%%%     The acknowledgements page         %%%%%%%%%%%%%%%%%%  
%%%%%%%%%%%%%%%%%%%%%%%%%%%%%%%%%%%%%%%%%%%%%%%%%%%%%%%%%%%%%%%%%%%%%%%%%%%
%\chapter*{Agradecimentos}
%\addcontentsline{toc}{chapter}{\numberline{}Acknowledgements}
%Thank to someone who prepared this template. Thank to someone who
%prepared this template. Thank to someone who prepared this template.
%Thank to someone who prepared this template. Thank to someone who
%prepared this template. Thank to someone who prepared this template.
%Thank to someone who prepared this template. Thank to someone who
%prepared this template. Thank to someone who prepared this
%template.Thank to someone who prepared this template. Thank to someone
%who prepared this template. Thank to someone who prepared this
%template. Thank to someone who prepared this template. Thank to
%someone who prepared this template. Thank to someone who prepared this
%template. Thank to someone who prepared this template. Thank to
%someone who prepared this template.

%%%%%%%%%%%%%%%%%%%%%%%%%%%%%%%%%%%%%%%%%%%%%%%%%%%%%%%%%%%%%%%%%%%%%%%%%%%
%%%%%%%%    tableofcontents, listoffigures and listoftables       %%%%%%%%%
%%%%%%%%        Command if you do not have  them                  %%%%%%%%%
%%%%%%%%%%%%%%%%%%%%%%%%%%%%%%%%%%%%%%%%%%%%%%%%%%%%%%%%%%%%%%%%%%%%%%%%%%%
\tableofcontents
\listoffigures
\listoftables
\clearpage


%%%%%%%%%%%%%%%%%%%%%%%%%%%%%%%%%%%%%%%%%%%%%%%%%%%%%%%%%%%%%%%%%%%%%%%%%%%
%%%%%%%%%%%%%%%%%%%%%%   END OF FRONT PAGE %%%%%%%%%%%%%%%%%%%%%%%%%%%%%%%%
%%%%%%%%%%%%%%%%%%%%%%%%%%%%%%%%%%%%%%%%%%%%%%%%%%%%%%%%%%%%%%%%%%%%%%%%%%%









